% 坐标投影：(a)点质量0.2@(0,0),0.3@(0,1),0.5@(1,1)，横轴边缘0.5delta_0+0.5delta_1，纵轴边缘0.2delta_0+0.8delta_1；(b)单位正方形面积测度，竖条B×R与正方形之交的面积等于lambda(B∩[0,1])。
\begin{tikzpicture}[x=1cm,y=1cm,>=Stealth,
  note/.style={font=\FigureNoteFont\fontsize{7.5}{9}\selectfont,text=FigureMuted},
  tl/.style={font=\fontsize{7.5}{9}\selectfont}]
\begin{scope}
  \draw[FigureGrid,dashed,line width=.5pt] (0.8,-0.75)--(0.8,2);\draw[FigureGrid,dashed,line width=.5pt] (2,-0.75)--(2,2);
  \draw[FigureGrid,dashed,line width=.5pt] (-0.75,0.8)--(0.8,0.8);\draw[FigureGrid,dashed,line width=.5pt] (-0.75,2)--(2,2);
  \foreach \x/\y/\m/\r in {0.8/0.8/0.2/2.2pt,0.8/2/0.3/2.7pt,2/2/0.5/3.4pt}{\fill[FigureBlue] (\x,\y) circle (\r);\node[tl,anchor=south west] at (\x+0.05,\y+0.03) {$\m$};}
  \draw[FigureStroke,line width=.6pt,->] (0.2,-0.75)--(2.7,-0.75) node[right,tl] {$x_1$};
  \draw[FigureStroke,line width=.6pt,->] (-0.75,0.2)--(-0.75,2.7) node[above,tl] {$x_2$};
  \foreach \x/\h/\l in {0.8/0.5/0,2/0.5/1}{\fill[FigureGreen] (\x-0.08,-0.75) rectangle (\x+0.08,-0.75+\h);\node[tl,below] at (\x,-0.75) {$\l$};\node[note,anchor=west] at (\x+0.1,-0.45) {$0.5$};}
  \foreach \y/\w/\l in {0.8/0.2/0,2/0.8/1}{\fill[FigureYellow] (-0.75,\y-0.08) rectangle (-0.75+\w,\y+0.08);\node[tl,left] at (-0.75,\y) {$\l$};}
  \node[note] at (-0.45,0.5) {$0.2$};\node[note] at (-0.35,2.3) {$0.8$};
  \node at (1.0,-1.55) {(a) 点质量的两个边缘};
\end{scope}
\begin{scope}[xshift=5.4cm]
  \fill[FigureBlueFill] (0,0) rectangle (2,2);
  \fill[FigureBlue!45] (0.6,0) rectangle (1.2,2);
  \draw[FigureBlue,line width=.6pt] (0,0) rectangle (2,2);
  \draw[FigureGreen,dashed,line width=.6pt] (0.6,-0.3)--(0.6,2.5);\draw[FigureGreen,dashed,line width=.6pt] (1.2,-0.3)--(1.2,2.5);
  \draw[FigureStroke,line width=.6pt,->] (-0.3,0)--(2.9,0) node[right,tl] {$x_1$};
  \draw[FigureStroke,line width=.6pt,->] (0,-0.3)--(0,2.6) node[above,tl] {$x_2$};
  \node[tl,below] at (2,0) {$1$};\node[tl,left] at (0,2) {$1$};
  \draw[FigureGreen,line width=2.4pt] (0.6,-0.6)--(1.2,-0.6);
  \node[tl] at (0.9,-0.9) {$B$};
  \node[note,anchor=west] at (2.15,1.0) {$B\times\mathbb R$};
  \node at (1.1,-1.55) {(b) 面积测度的边缘};
\end{scope}
\end{tikzpicture}
